En bikonvex lins av glas med brytningsindex $n = 1{,}50$ har krökningsradierna $r_1 = 20\,\mathrm{cm}$ (första ytan) och $r_2 = -30\,\mathrm{cm}$ (andra ytan).
\begin{enumerate}
  \item Beräkna linsens brytningsstyrka och fokallängd i luft.
  \item Linsen sänks ned i vatten ($n = 1{,}33$). Använd sambandet för avbildning genom en sfärisk gränsyta,
  \[
    \frac{n}{s} + \frac{n'}{s'} = \frac{n' - n}{r} = P,
  \]
  på var och en av linsens två ytor för att beräkna linsens brytningsstyrka i vatten.
  \item Var hamnar bilden av ett mycket avlägset objekt när linsen är i vatten? Jämför med fokallängden i luft.
  \item Använd resultatet för att förklara varför allt ser suddigt ut när man öppnar ögonen under vatten utan simglasögon.
\end{enumerate}

\emph{Ledning:} För en tunn lins gäller att bilden från den första ytan blir objekt för den andra, med $s_2 = -s_1'$. Omgivningen är nu vatten på båda sidor.
